\chapter{放射危害评价}

\section{正常运行条件下的放射危害评价}

根据第2章2.6工程分析章节和第3章第3.2正常运行状态下的辐射源项分析章节可知，本项目在建成后，正常运行条件下，工作人员和公众受到的辐射危害情况见表~\ref{tab:ch6-1-t1-a18f249f}。


\begin{table}[H]
\centering
\caption{正常情况工作人员和公众受到辐射危害的情况分析}
\label{tab:ch6-1-t1-a18f249f}
\begin{tabular}{|p{2.5cm}|p{2.5cm}|p{6cm}|p{2.5cm}|}
\hline
\multicolumn{2}{|c|}{项目} & 工作人员 & 公众 \\
\hline
放射治疗项目 & 后装治疗机 & 
1) 在机房外控制室操作 \newline
2) 在治疗前工作人员在机房内摆位 & 
在机房外周围停留 \\
\hline
\end{tabular}
\end{table}


\subsection{放射治疗项目正常运行条件下的放射危害评价}

本项目后装治疗机机房外放射工作人员、公众成员所接受的年附加有效剂量估算结果下表~\ref{tab:ch6-1-t2-9c5297f7}。


\begin{table}[H]
\centering
\caption{本项目机房外人员年剂量估算结果一览表}
\label{tab:ch6-1-t2-9c5297f7}
% 只改动这里：用 p 列固定宽度，其余完全和你原版一样
\begin{tabular}{|m{1cm}|m{1cm}|m{2cm}|m{1cm}|m{1.8cm}|m{1cm}|m{1cm}|m{1.6cm}|m{1.6cm}|}
\hline
\multirow{2}{=}{机房名称}
 & \multicolumn{3}{c|}{关注区域}
 & \multicolumn{5}{c|}{年有效剂量计算结果} \\
\cline{2-9}
 & 关注点
 & 对应区域
 & 涉及成员
 & 剂量率\ding{172} ($\mu$Sv/h)
 & 居留因子$T$\ding{172}
 & 使用因子$U$
 & 年工作时间$t$（h）
 & 年剂量\ding{174} (mSv) \\
\hline
\multirow{6}{=}{后装治疗机房}
 & a & 准备间      & 公众    & 6.33E-03 & 1/2  & 1 & 254 & 8.04E-04 \\
\cline{2-9}
 & b & 狭窄过道    & 公众    & 3.89E-02 & 1/40 & 1 & 254 & 2.47E-04 \\
\cline{2-9}
 & c & 直线加速器2 & 工作人员& 2.49E-02 & 1/8  & 1 & 254 & 7.91E-04 \\
\cline{2-9}
 & d & 控制室      & 公众    & 5.08E-02 & 1    & 1 & 254 & 1.29E-02 \\
\cline{2-9}
 & e & 屋顶平台    & 公众    & 8.61E-03 & 1/40 & 1 & 254 & 5.47E-05 \\
\cline{2-9}
 & f & 防护门外    & 公众    & 3.13E-01 & 1/8  & 1 & 254 & 9.94E-03 \\
\hline
\multicolumn{9}{|m{\dimexpr\textwidth-2\tabcolsep-2\arrayrulewidth}|}{
注：1）居留因子取值依据见表~\ref{tab:ch4-dose-control}；\newline
2）剂量率由屏蔽厚度估算出的剂量率（取预估结果最大值，见表~\ref{tab:ch4-verify}）。\newline
3）年工作时间见上文2.6章节。\newline
4）年剂量（mSv）=剂量率（$\mu$Sv/h）$\times$使用因子$\times$居留因子$\times$年工作负荷（h/年）/1000。
} \\
\hline
\end{tabular}
\end{table}


\subsubsection{放射工作人员预计接受的剂量估算}

根据建设单位提供的工作量，假设建设单位按照设计屏蔽情况进行施工，其各机房外人员可能受到的年有效剂量如表~\ref{tab:ch6-1-t2-9c5297f7}所示。

\textbf{（1）后装治疗机按照极不利的条件下估算：}对于后装治疗技师需考虑摆位过程中受贮源器表面泄漏辐射所致有效剂量，根据标准 WS 262-2017，对于贮源器表面1m处泄漏辐射所致周围剂量当量率不超过 $5 \mu \mathrm { S v / h }$，后装机每年治疗 1250 人次，每人次摆位平均时间为 1min。摆位由 2 名技师进行，则每位摆位工作人员年剂量$= 5 \mu \mathrm { S v } / \mathrm { h } { \times } 1 2 5 0$人/年×1min÷（2 人 $\times 6 0 \mathrm { m i n / h } { \times } 1 0 0 0 \mu \mathrm { S v / m S v }$） ${ \approx } 0 . 0 5 2 \mathrm { m S v }$。则后装治疗技师工作人员可能受到的年附加剂量为摆位时受到的附加剂量（ $0 . 0 5 2 \mathrm { m S v }$），叠加在后装治疗机房外曝光时所受到的附加剂量之和约为 $0 . 0 6 5 \mathrm { m S v }$。

综上所述，正常情况下，放射治疗项目放射工作人员年有效剂量小于国家限值并小于建设单位为本项目设置的管理目标（5.0mSv）。

\subsubsection{公众预计接受的辐射剂量估算}

对于后装治疗机机房周围公众曝光时所受到的剂量为 $9 . 9 4 \mathrm { E } { - } 0 3 \mathrm { m S v }$，小于国家限值并小于建设单位为本项目设置的管理目标（ $0 . 2 5 \mathrm { m S v }$）。

\subsubsection{本项工作人员预计接受的辐射剂量与原项目累积估算}

由以上内容可知，本项目职业人员附加年有效剂量最大为 $0 . 0 6 5 \mathrm { m S v }$，保守叠加原有工作负荷后（见附件 9-4，取近一年最大年有效剂量人员（（放射工作人员））的剂量进行叠加），年有效剂量最大为 $0 . 2 \mathrm { m S v }$，低于年有效剂量约束值 5mSv，工作人员附加年有效剂量最大为 $0 . 8 5 \mathrm { m S v }$，小于国家限值并小于建设单位为本项目设置的管理目标（5mSv）。

综上所述，按照本项目放射治疗用房关注区域按设计屏蔽厚度计算的周围剂量当量率理论估算值计算得出，相应人员预计接受和年剂量均满足标准要求及管理目标值。

\section{异常和事故情况下放射危害评价}

\textbf{对于后装治疗项目作出如下相应的预防措施：}

1）操作系统发生故障，如后装治疗机发生卡源，到达预置治疗时间，放射源无法自动收回，导致患者受超剂量照射。

预防措施：按规章制度定期对各个联锁装置进行检查，发现故障及时清除，严禁在门－机装置失效的情况下违规操作；通过装置故障报警系统及时发现故障，及时修复；通过纵深防御以减少由于某个联锁失效或在某个联锁失效期间产生的辐射，定期对安全联锁装置进行自检，禁止在失灵情况下运行。

2）换装放射源后的废源因管理不善发生被盗、丢失等事故，可能造成对公众的意外照射和环境放射性污染。

预防措施：建设单位拟委托厂家负责将废旧放射源返回原厂家处理，废源不在建设单位存储。

3）更换放射源或维修时，放射源意外破损时，治疗装置和机房还可能会存在放射性污染。

预防措施：建设单位拟委托厂家负责源的更换，不自主更换源，如发生放射源意外破损时，存储在应急储存设施内。

4）人员误入或滞留后装治疗机房造成的误照射。

预防措施：建设单位拟安装摄像监控设施，撤离机房时应清点人数，按程序对机房全视角搜寻；用对讲机呼叫，声灯报警。运行时发现有人员滞留室内，由控制室紧急按下停机开关；如机房内有工作人员未及时出来，可通过随身携带的仪器报警，立即按下机架上的紧急停机开关，可将辐射危害的严重程度降至最低限度；机房内摄像监控装置实时观察机房内动态，定期检查摄像监控装置的有效性。

5）治疗计划失误，如制订治疗计划时，弄错患者或剂量或分次剂量与处方数值严重不符。

预防措施：放射工作人员须加强专业知识学习，加强职业技能培训，增强责任感，严格遵守操作规程和规章制度；管理人员应强化管理，落实安全责任制，经常督促检查。

6）屏蔽设施损坏或老化达不到放射防护要求，使射线泄漏到控制室，可能对控制室长期驻留的医护人员造成小剂量的长期照射。

预防措施：建设单位拟规定对机房外周围剂量当量率每月定期进行检测。

潜在照射是指在正常情况下有一定把握预期不会发生，但可能会因为事故或某种偶然性质的事件或时间序列（包括设备故障和操作错误）所引起的照射。通常是指有一定发生概率而不一定发生的照射。潜在照射是不能预先完全计划或知道的照射，常是由于出现意外情况或事故的照射。在潜在照射发生前的防护对策，主要是预防和缓减，减少概率的措施称为预防，降低剂量大小或严重程度的措施称为缓减。通常潜在照射的发生概率和剂量大小都可在一定程度上加以控制。从第3章3.3异常或事故状态下的辐射源项章节介绍可知，建设项目异常或事故状态下的辐射源项为γ射线所产生的轫致辐射，建设单位对建设项目潜在照射情况进行预防，符合相关要求。

\section{小结}

本项目在正常运行情况下，相关人员的年有效剂量可控制在建设单位的管理目标值范围内，并计划通过个人剂量计定期监控放射工作人员的年有效剂量。建设单位对异常或事故情况受照人员进行分析，并提出了相应的预防措施，以降低其发生的概率，符合相关要求。
